%\chapter{The Optimum Receiver}

%\subsection{Problem}
\begin{definition}


A right angled triangle looks like Fig. \ref{ch1_right}.
\begin{figure}[!h]
\begin{center}
	
\includegraphics[width=\columnwidth]{./figs/ch1_right_triangle}
\vspace*{-10cm}
\end{center}
\caption{Right Angled Triangle}
\label{ch1_right}	
\end{figure}
%\vspace*{-10cm}
with angles $\angle A,\angle B$ and $\angle C$ and sides $a, b$ and $c$.  The unique feature of this triangle is $\angle C$ which is defined to be $90^{\degree}$.
\end{definition}
\begin{definition}
	For simplicity, let the greek letter $\theta = \angle B$.  We have the following definitions.
\begin{equation}
\label{ch1_trig_defs}
\begin{matrix}
	\sin \theta = \frac{a}{c} & 	\cos \theta = \frac{b}{c} \\
	\tan \theta = \frac{b}{a} & \cot \theta = \frac{1}{\tan \theta} \\
	\csc \theta = \frac{1}{\sin \theta} & \sec \theta = \frac{1}{\cos \theta}
	\end{matrix}
\end{equation}
\end{definition}

\subsection{Sum of Angles}
\begin{figure}[!h]
	\begin{center}
		
		\includegraphics[width=\columnwidth]{./figs/ch1_parallel_triangle}
		\vspace*{-10cm}
	\end{center}
	\caption{Sum of angles of a triangle}
	\label{ch1_parallel_triangle}	
\end{figure}

\begin{definition}
	In Fig. \ref{ch1_parallel_triangle}, the sum of all the angles on the top or bottom side of the straight line $XY$ is $180^{\degree}$.
\end{definition}

\begin{definition}
In Fig. \ref{ch1_parallel_triangle}, the straight line making an angle of $90^{\degree}$ to the side $AC$ is said to be parallel to the side $BC$. Note there is an angle at $A$ that is equal to $\theta$.  This is one property of parallel lines.  Thus, $\angle YAZ = 90^{\degree}$.
\end{definition}

\begin{problem}
	Show that $\angle VAZ = 90^{\degree} - \theta$
\end{problem}		
	\proof Considering the line $XAZ$,
	\begin{align}
	\theta + 90^{\degree} + \angle VAZ &= 180^{\degree} \\
	\Rightarrow  \angle VAZ =  90^{\degree} - \theta
	\end{align}

\begin{problem}
	\label{ch1_compl_angle}
	Show that $\angle BAC = 90^{\degree} - \theta$.
\end{problem}	
	\proof Consider the line $VAB$ and and use the approach in the previous problem.  Note that this implies that $\angle VAZ = \angle BAC$.  Such angles are known as vertically opposite angles. 
	 
\begin{problem}
Sum of the angles of a triangle is equal to $180^{\degree}$
\end{problem}
\subsection{Budhayana Theorem}



\begin{figure}[!h]
	\begin{center}
		
		\includegraphics[width=\columnwidth]{./figs/ch1_budh_triangle}
		\vspace*{-10cm}
	\end{center}
	\caption{Budhayana Theorem}
	\label{ch1_budh_triangle}	
\end{figure}

\begin{problem}
	Using Fig. \ref{ch1_right}, show that
	\begin{equation}
	\cos \theta = \sin \brak{90^{\degree} - \theta}
	\end{equation}
\end{problem}
\proof From Problem \ref{ch1_compl_angle} and  \eqref{ch1_trig_defs}
%
\begin{equation}
	\cos \brak{90^{\degree}-\theta} = \frac{b}{c} = \sin \theta
\end{equation}
%
\begin{problem}
Using Fig. \ref{ch1_budh_triangle}, show that 
%
\begin{equation}
\label{ch1_budh_basic}
c = a \cos \theta + b \sin \theta
\end{equation}
%
\end{problem}
\proof We observe that
%
\begin{align}
BD &= a \cos \theta \\
AD &= b \cos\brak{90-\theta} = b \sin \theta \quad \brak{\text{From} \quad \eqref{ch1_compl_angle}}
\end{align}
%
Thus,
\begin{equation}
BD + AD = c = a \cos \theta + b \sin \theta
\end{equation}
\begin{problem}
From \eqref{ch1_budh_basic}, show that
%
\begin{equation}
\sin ^2 \theta + \cos ^2 \theta = 1
\end{equation}

\end{problem}
%
\proof Dividing both sides of \eqref{ch1_budh_basic} by $c$, 
\begin{align}
1 &= \frac{a}{c}\cos\theta + \frac{b}{c}\sin\theta\\
\Rightarrow &\sin ^2 \theta + \cos ^2 \theta = 1 \quad \brak{\text{from} \quad \eqref{ch1_trig_defs}}
\end{align}

\begin{problem}
	Using \eqref{ch1_budh_basic}, show that
	\begin{equation}
	\label{ch1_budhayana_them}
	c^2 = a^2 + b^2
	\end{equation}
	\eqref{ch1_budhayana_them} is known as the Budhayana theorem.  It is also known as the Pythagoras theorem.
\end{problem}
\proof From \eqref{ch1_budh_basic},
\begin{align}
c &= a\frac{a}{c} + b \frac{b}{c} \quad \brak{\text{from} \quad \eqref{ch1_trig_defs}}\\
\Rightarrow c^2 &= a^2 + b^2
\end{align}


